\chapter{جمع‌بندی و مسیرهای آینده}
\label{ch:conclusions}

\section{پاسخ به پرسش پژوهش}
این پایان‌نامه مسئله پیش‌بینی بیش‌برازش فاجعه‌بار در آموزش هم‌آورد سریع را از دیدگاه هندسه تابع هزینه بررسی کرد. نقطه آغاز آن، شکاف میان درستی زیاد در برابر حمله تک‌گامی و درستی اندک در برابر حمله چندگامی بود. این شکاف نشان می‌دهد یک مولد نمونه آموزشی ممکن است به‌تدریج از حل مسئله درونی بازبماند، درحالی‌که فرایند آموزش همچنان روی خروجی آن پیشرفت ظاهری نشان می‌دهد.

بیش‌برازش اپسیلون چارچوبی برای توضیح بخشی از این رفتار فراهم کرد: یادگیری روی اندازه‌ها و جهت‌های محدود آشفتگی می‌تواند نقاطی با هزینه کم در انتهای حرکت بسازد و نقاط آسیب‌پذیر میان آن‌ها را پوشش ندهد. شواهد تصویری مقاله از این تفسیر پشتیبانی می‌کنند، ولی این پدیده به‌عنوان علت کافی و یگانه همه موارد بیش‌برازش فاجعه‌بار اثبات نشده است.

برای پایش، شاخص \pertalign{} با مقایسه جهت شیب در آغاز و پایان حمله تعریف شد. تحلیل محلی نشان داد افت آن با بخشی از تغییر شیب که بر شیب اولیه عمود است ارتباط دارد. ازاین‌رو، این شاخص نشانه‌ای از نامعتبرشدن تقریب خطی در مسیر حمله فراهم می‌کند. اطلاعات لازم عمدتاً در همان دو مرحله تولید نمونه و به‌روزرسانی مدل حضور دارد؛ بنابراین، نیاز به ساخت هسین کامل یا اجرای منظم حمله چندگامی برای هر دسته کاهش می‌یابد.

\section{از پایش تا اصلاح آموزش}
در روش \sora{}، اختلاف دو شیب برای تقریب خمیدگی جهتی به کار رفت و یک قانون عملی برای کران بالای اندازه گام ساخته شد. میانگین‌گیری نمایی، تخمین را در برابر تغییر دسته‌ها هموار می‌کند و نمونه‌برداری مؤلفه‌ای اندازه گام، تنوع را حفظ می‌کند. این ترکیب، پیش‌بینی را به یک سازوکار اصلاحی پیوند می‌دهد: مدل تنها درباره افت هندسی هشدار نمی‌گیرد، بلکه توزیع ورودی آموزشی نیز تغییر می‌کند.

تمایز میان استدلال نظری و الگوریتم عملی در این نتیجه اهمیت دارد. حل دقیق یک مدل درجه دوم روی یک جهت ثابت، مبنای طراحی است. ثابت‌کردن صورت کسر، استفاده از تخمین دسته‌های گذشته، گام‌های مؤلفه‌ای و برش دامنه پیکسل، انتخاب‌های عملی‌اند. موفقیت تجربی این انتخاب‌ها، تضمین بهینگی سراسری مسئله آموزش یا حل دقیق بیشینه‌سازی درونی را نتیجه نمی‌دهد.

\section{برداشت از نتایج}
نتایج مقاله مشترک، روی مجموعه‌ای از دادگان طبیعی و زیست‌پزشکی و معماری‌های گوناگون، نشان می‌دهد \sora{} در بسیاری از تنظیم‌ها بده‌بستان مطلوبی میان درستی پاک، درستی هم‌آورد و هزینه ایجاد می‌کند \cite{sora2026}. اهمیت دادگان زیست‌پزشکی آن است که برخی روش‌های موفق در معیارهای رایج، در این تنظیم‌ها با افت جدی روبه‌رو می‌شوند. بنابراین، محدودکردن ارزیابی به یک معیار تصویری می‌تواند تفاوت روش‌ها را پنهان کند.

بااین‌حال، جدول‌های تفصیلی مقاله نشان می‌دهند برتری عددی در همه معیارها و همه ترکیب‌ها برقرار نیست. در نتیجه، برداشت مناسب، پایداری و رقابت‌پذیری گسترده در تنظیم‌های بررسی‌شده است. همچنین، اعداد متن اصلی، حذف مؤلفه‌ها و پیوست باید در زمینه آزمایش خود خوانده شوند؛ اختلاف آن‌ها را نمی‌توان بدون گزارش خام و مشخصات نقطه کنترل به یک نتیجه واحد تبدیل کرد.

\section{محدودیت‌های باقی‌مانده}
نخست، هشدار زودهنگام هنوز نیازمند ارزیابی آماری مستقل است. افت پیش از فروپاشی در یک نمودار، وجود یک آستانه عمومی با نرخ هشدار کاذب کنترل‌شده را اثبات نمی‌کند. به‌ویژه، تغییرات طبیعی شیب در آغاز آموزش یا تغییر نرخ یادگیری می‌تواند بر شاخص اثر بگذارد. برای استفاده عملی، باید مشخص شود چه افتی معنادار است و چه مدت تداوم لازم دارد.

دوم، شاخص یک خلاصه جهت‌وابسته است. ممکن است شیب در راستای خود به‌شدت تغییر بزرگی دهد ولی هم‌راستا باقی بماند؛ یا خمیدگی در جهت‌های بررسی‌نشده زیاد باشد. ترکیب \pertalign{} با نسبت بزرگی شیب، هزینه و یک ارزیابی چندگامی کم‌تواتر می‌تواند تصویری کامل‌تر بدهد، اما هزینه و سود چنین ترکیبی باید اندازه‌گیری شود.

سوم، استفاده از آمار دسته‌های گذشته به شباهت توزیعی و تغییرات نسبتاً آرام مدل متکی است. در دسته‌های نامتوازن، یادگیری پیوسته یا تغییر ناگهانی توزیع، یک ضریب جهانی ممکن است رفتار همه نمونه‌ها را به‌خوبی بازنمایی نکند. همچنین، دامنه آشفتگی آموزش \sora{} با کره تهدید ارزیابی یکسان نیست؛ هر مقایسه آینده باید این نکته را آشکارا گزارش کند.

\section{پیشنهادهای ادامه پژوهش}
یک مسیر مشخص، تبدیل شاخص هندسی به پیش‌بین واسنجی‌شده است. باید رویداد بیش‌برازش با قاعده‌ای مستقل از خود شاخص تعریف شود، آستانه روی داده اعتبارسنجی تعیین شود و سپس نرخ هشدار درست، هشدار کاذب و زمان پیش‌آگاهی روی اجراهای کنارگذاشته‌شده گزارش شود. چند بذر تصادفی و چند معماری برای جداکردن یک الگوی پایدار از نوسان یک اجرا ضروری است.

مسیر دوم، مطالعه کنترل در سطح نمونه یا گروه است. به‌جای یک میانگین جهانی، می‌توان توزیع شاخص را در هر دسته یا در دسته‌های معنایی دنبال کرد. این کار ممکن است برای داده نامتوازن سودمند باشد، اما افزایش حافظه، نوسان و وابستگی به شناسه نمونه را به همراه دارد. مقایسه منصفانه باید نشان دهد آیا سود ستبری، هزینه این پیچیدگی را توجیه می‌کند.

مسیر سوم، مطالعه جداگانه نقش تنوع، دامنه آشفتگی و تطبیق هندسی است. مقایسه روش‌ها با دامنه آموزش یکسان و نیز با بودجه محاسباتی یکسان کمک می‌کند سهم هریک روشن شود. ارزیابی در برابر بودجه‌های دیده‌نشده، حملات متفاوت و جابه‌جایی توزیع نیز می‌تواند حدود تعمیم را دقیق‌تر مشخص کند.

در نهایت، پیوند با کران‌های نظری ستبری و با شبکه‌های غیرهموار همچنان باز است. تحلیل فعلی توضیحی محلی برای یک سنجه و یک قانون کنترل فراهم می‌کند. توسعه آن به تضمین‌هایی درباره درستی دسته‌بندی یا پایداری کل فرایند یادگیری، به فرض‌های تازه و تحلیل مستقل نیاز دارد. ارزش عملی مسیر حاضر، نشان‌دادن امکان استفاده از اطلاعات هندسی موجود در آموزش برای کاهش شکست یک تقریب سریع است.
